This book is written for advanced undergraduates, graduate students, and
practising engineers who want a unified, self-contained account of why reactors
are (or are not) stable, and what happens physically when stability is lost. It
assumes a working knowledge of multivariable calculus, ordinary differential
equations, and undergraduate physics. No prior course in nuclear engineering is
strictly required: the necessary reactor physics is built up from the neutron
balance in Chapters~2 and~3.

\section*{Organisation}

The book is organised into five parts and a set of appendices.

\textbf{Part I (Foundations of Reactor Dynamics)} establishes the language:
fission, the multiplication factor and reactivity, the point reactor kinetics
equations and the indispensable role of delayed neutrons, and the decay-heat
source that persists long after the chain reaction stops.

\textbf{Part II (Heat Transport and Reactivity Feedback)} follows the energy from
its birth in the fuel to its removal by the coolant, and then introduces the
feedback coefficients that close the loop between temperature and reactivity.

\textbf{Part III (Stability Theory)} treats stability quantitatively, first by
linearisation and frequency-domain methods, then through the nonlinear physics
of power excursions, and finally through the slow, spatial instabilities driven
by xenon-135.

\textbf{Part IV (Historical Lessons)} applies the theory to real machines. After
a survey of earlier events, it examines the RBMK-1000 design in detail, then
reconstructs and analyses the Chernobyl accident.

\textbf{Part V (Modern Practice and Synthesis)} describes how the lessons are
embodied in contemporary designs that aim for inherent and passive safety, and
demonstrates the dynamics computationally before drawing the threads together.

\section*{A note on conventions}

Reactivity is expressed both as an absolute fraction $\rho$ (units of $\Delta
k/k$, often quoted in per-cent-mille, pcm, where $1\unit{pcm}=10^{-5}$) and in
\emph{dollars}, $\rho/\beta$, where $\beta$ is the total delayed-neutron
fraction. One dollar of reactivity is the boundary of prompt criticality and is
the natural yardstick for transient severity. Temperatures are in kelvin or
degrees Celsius as convenient. SI units are used throughout. Numerical examples
use validated delayed-neutron data for $^{235}$U; the full data set appears in
Appendix~A.

\section*{Companion software}

Several figures and worked transients in this book were produced with a small,
open-source point-reactor-kinetics toolkit written by the author. The code,
which couples validated kinetics data to a stiff ODE solver with temperature
feedback, is referenced in Chapter~16 and is freely available. Readers are
encouraged to reproduce and extend the examples.

\section*{Scope and disclaimer}

This is an educational and historical work. It deliberately discusses reactor
behaviour at the level of physical principle and published, openly available
engineering knowledge. It is \emph{not} an operating procedure, a design
specification, or a safety-analysis report, and nothing in it should be used as
such.
